\documentclass{article}
\usepackage{amsmath, amssymb, amsfonts}
\usepackage{geometry}
\usepackage{booktabs}
\usepackage{xcolor}
\geometry{margin=1in}
\title{UV Check Mesh Audit}
\author{Rune}
\date{2026-05-11}

\begin{document}
\maketitle

\section{Cube: 5 of 6 Faces Have Reversed Winding}

Face normals computed from index winding (edge1 $\times$ edge2):

\begin{table}[h]
\centering
\begin{tabular}{lccc}
\toprule
\textbf{Face} & \textbf{Winding normal} & \textbf{Outward} & \textbf{Correct?} \\
\midrule
Front (+Z) & $-Z$ & $+Z$ & \color{red}\textbf{WRONG} \\
Back  (-Z) & $+Z$ & $-Z$ & \color{red}\textbf{WRONG} \\
Right (+X) & $-X$ & $+X$ & \color{red}\textbf{WRONG} \\
Left  (-X) & $-X$ & $-X$ & \color{green}{OK} \\
Top   (+Y) & $-Y$ & $+Y$ & \color{red}\textbf{WRONG} \\
Bottom(-Y) & $+Y$ & $-Y$ & \color{red}\textbf{WRONG} \\
\bottomrule
\end{tabular}
\end{table}

Only the Left face (indices in sequential order \texttt{12,13,14, 12,14,15}) has correct
CCW-from-outside winding. All other faces swap the 2nd and 3rd vertices in each triangle,
reversing the winding.

With \texttt{VK\_CULL\_MODE\_NONE} this doesn't prevent rendering, but it means the
rasterizer treats these triangles as \emph{back-facing}. Since the glTF loader and Blender
produce CCW-wound meshes, the procedural cube is inconsistent.

\textbf{Fix}: swap indices 2 and 3 in each triangle to match CCW convention:

\begin{verbatim}
  idxs =
    [ -- Front:  was 0,2,1  0,3,2
      0, 1, 2, 0, 2, 3,
      -- Back:   was 4,6,5  4,7,6
      4, 5, 6, 4, 6, 7,
      -- Right:  was 8,10,9  8,11,10
      8, 9, 10, 8, 10, 11,
      -- Left:   already correct
      12, 13, 14, 12, 14, 15,
      -- Top:    was 16,17,18  16,18,19
      16, 19, 18, 16, 18, 17,
      -- Bottom: was 20,22,21  20,23,22
      20, 21, 22, 20, 22, 23
    ]
\end{verbatim}

\section{Cube: ``Upside-down on all faces'' --- Texture File Convention}

The UVs in the code are mathematically correct for Vulkan convention
($v{=}0$ at physical top of each face, $v{=}1$ at bottom). I verified each face
by tracing camera $\to$ view $\to$ clip $\to$ NDC $\to$ texture coordinates.

The 180° rotation on ALL faces (confirmed by vision model) most likely means
\texttt{data/textures/uv\_checker.png} has its content oriented for OpenGL convention
($v{=}0$ at bottom). When sampled with Vulkan convention ($v{=}0$ at top),
every face shows the texture 180° rotated.

\textbf{Evidence}: the glTF test cube works correctly, but it uses
\emph{different textures} (embedded in the glTF file), not \texttt{uv\_checker.png}.
So the ``glTF works, procedural doesn't'' comparison is texture-vs-texture,
not UV-vs-UV.

\textbf{Fix}: verify \texttt{uv\_checker.png} orientation in an image viewer.
If text reads correctly top-to-bottom, the texture is fine and the issue is elsewhere.
If text is stored bottom-to-top (OpenGL convention), regenerate the texture or
flip V during loading.

\section{Sphere: U Coordinate Inverted}

\begin{verbatim}
u = realToFrac (1.0 - fromIntegral lonIdx / fromIntegral lonCount)
\end{verbatim}

At $\phi = 0$ (direction $+X$ at equator): $u = 1.0$ (right edge of texture).
At $\phi = \pi$ (direction $-X$): $u = 0.5$ (center).

Standard convention maps $\phi = 0 \to u = 0$. The current code mirrors the
texture horizontally.

\textbf{Fix}:
\begin{verbatim}
u = realToFrac (fromIntegral lonIdx / fromIntegral lonCount)
\end{verbatim}

\section{Sphere: Tangent Needs Updating}

The tangent is correct for the current (inverted) U mapping:
\begin{verbatim}
tx = realToFrac (sin phi)
tz = realToFrac (-cos phi)
\end{verbatim}

This points in the direction of \emph{decreasing} $\phi$ (= increasing U with the
inverted mapping). After fixing U, the tangent must point in the direction of
\emph{increasing} $\phi$:

\textbf{Fix}:
\begin{verbatim}
tx = realToFrac (-sin phi)
tz = realToFrac (cos phi)
\end{verbatim}

Or equivalently, negate the tangent: \texttt{V4 (-tx) (-ty) (-tz) tw}.

\section{Sphere: Reversed Winding (Same as Cube)}

Triangle order \texttt{base, nextBase, nextBase+1} produces inward-pointing
face normals (verified numerically at equator). The fix is to reverse the
triangle order to \texttt{base, base+1, nextBase, base, nextBase, nextBase+1}:

\begin{verbatim}
    in [ fromIntegral base
       , fromIntegral (base + 1)
       , fromIntegral nextBase
       , fromIntegral base
       , fromIntegral nextBase
       , fromIntegral (nextBase + 1)
       ]
\end{verbatim}

\section{Sphere: ``Half Reflective, Half Non-Reflective''}

This is most likely caused by the \textbf{environment cubemap orientation} combined
with the \textbf{reversed winding}, not the sphere geometry itself.

With inward-pointing face normals, the rasterizer marks these fragments as
back-facing. While \texttt{CULL\_MODE\_NONE} still renders them, some GPU
pipelines may handle back-facing fragments differently in subtle ways (e.g.,
helper invocations, derivative computations for cubemap LOD selection).

Additionally, the cubemap faces loaded from \texttt{data/hdri/env\_test/} may be
in OpenGL convention (see \texttt{SHADER\_PIPELINE\_MATH.tex} Issue 1). If some
faces are horizontally flipped, reflection sampling returns wrong texels for
half the sphere directions, producing the visible hard seam.

\textbf{Recommended debugging}: enable debug mode 12 (raw skybox, all pixels)
and verify the cubemap renders correctly before investigating sphere reflections.

\section{Summary of Fixes}

\begin{table}[h]
\centering
\begin{tabular}{llp{5cm}}
\toprule
\textbf{\#} & \textbf{Priority} & \textbf{Fix} \\
\midrule
1 & High & Cube: fix 5 reversed-winding faces in \texttt{idxs} \\
2 & High & Sphere: fix U to \texttt{lonIdx/lonCount} (remove inversion) \\
3 & High & Sphere: fix tangent direction for corrected U \\
4 & High & Sphere: reverse triangle winding \\
5 & Medium & Verify \texttt{uv\_checker.png} orientation in image viewer \\
6 & Low & Cubemap face orientation (OpenGL vs Vulkan, per existing doc) \\
\bottomrule
\end{tabular}
\end{table}

\end{document}
